\section{Related Work}

\paragraph{CLIP and compositional failures.}
CLIP learns transferable image representations from natural-language supervision and remains a common frozen backbone for retrieval and zero-shot transfer~\cite{radford2021clip,ilharco2021openclip}. Several benchmarks show that strong image-text encoders can still behave like bags of words or exploit benchmark artifacts rather than bind objects, attributes, and relations faithfully~\cite{yuksekgonul2022bags,hsieh2023sugarcrepe,thrush2022winoground}. Our experiments use ARO- and SugarCrepe-style relation suites as staged counterfactual checks, but we also include Winoground and Flickr30k follow-ups because benchmark-specific gains do not automatically imply broad compositional generalization.

The staged suites are useful because they create controlled caption alternatives: many words remain fixed while an order, object, attribute, or relation changes. That format makes the relation failure easy to measure with a dual encoder. It is also limited. A model can learn to separate benchmark-style edits without acquiring the full visual grounding needed for natural paired-image binding. Winoground was included for exactly this reason. It asks for paired discrimination across two images and two captions, which is closer to a binding test than a single counterfactual ranking margin. The negative Winoground result in Section~4 therefore narrows the paper's claim rather than contradicting it.

Recent benchmark and prompting studies point to the same caution from different directions. Multi-image relation tests, causal-order probes, and visual causal reasoning suites show that modern vision-language models can struggle when the correct answer depends on inter-image structure, event order, or counterfactual intervention rather than object presence alone~\cite{zong2024mirb,komanduri2025causalvlbench,weng2026causalorder}. Compositional chain-of-thought prompting can elicit more explicit decompositions from large multimodal models~\cite{mitra2024ccot}, but RRB targets the lower-level dual-encoder setting where inference remains a single dot product between independently encoded image and text embeddings.

\paragraph{Structured image-text matching.}
Graph and structure-aware image-text retrieval methods explicitly represent objects and relations, often through scene graphs, relationship-enhanced semantic graphs, triple-level relation reasoning, dynamic multi-relational graphs, or multi-level semantic matching~\cite{wang2019crossmodal,liu2024resg,ma2026tsgr2,liu2024semscene,song2025cmrgr,chen2024spatialvlm}. These methods motivate the relation failure mode, but they differ from RRB's constraint: inference remains a normal dual encoder with no detector, parser, cross-encoder, graph extractor, or reranker. RRB uses structured counterfactuals as training supervision while keeping the runtime embedding interface unchanged.

This distinction is practical as well as conceptual. A cross-encoder or graph matcher can inspect a candidate image-caption pair jointly, but a CLIP-style retrieval index usually relies on precomputed embeddings and dot products. Changing that interface changes storage, indexing, latency, and downstream compatibility. RRB deliberately avoids that route. The method can be viewed as asking how much relation sensitivity can be added while still exporting one image vector and one text vector. This constraint makes the method weaker than explicit reasoning systems in some settings, but it makes the preservation problem more precise.

A related line uses structured knowledge or explicit visual structure to improve multimodal reasoning and pretraining. Examples include 3D scene-graph guided pretraining, unified multimodal knowledge graph protocols, annotation-free multimodal knowledge-graph construction, graph-conditioned reasoning with large language models, and structural-cue enhanced CLIP training~\cite{liu20243dscenegraph,gong2024uknow,liu2025valik,sun2026mario,ruan2026structxlip}. These papers strengthen the motivation for relational structure, but most either introduce extra structured inputs, graph construction, or model-scale reasoning machinery. RRB is deliberately smaller: it asks whether relation supervision can be absorbed into a frozen OpenCLIP-style embedding interface without changing retrieval-time inputs.

There is also a more distant family of latent-interface work that transfers internal states directly, such as KV-cache alignment or cache-to-cache communication between language models~\cite{dery2026latentcache,fu2026cachetocache}. RRB is not a communication protocol and does not pass caches between models, but this nearby literature is useful context for treating internal representation interfaces as design objects rather than only as hidden implementation details.

\paragraph{Relation-aware CLIP adaptation.}
Recent relation-aware CLIP work, including RACA-CLIP-style alignment efforts, targets compositional relation sensitivity directly. The local RACA-CLIP manuscript available in this project is under-review with incomplete bibliographic metadata, so we treat it as conceptual prior art rather than assign unsupported venue metadata~\cite{raca2026local}. The distinction is that RRB focuses on preserving the frozen model's retrieval and zero-shot anchor through a bounded residual path. The novelty is not ``using relations'' alone; it is the residualized tradeoff analysis plus the evidence that hard replacement bottlenecks are unsafe for broad alignment.

That positioning is intentionally modest. Relation supervision, counterfactual captions, and CLIP adaptation are active areas, so the contribution cannot rest on the existence of relational data. The evidence in this project is instead diagnostic: the same broad recipe can yield very different guardrail behavior depending on whether the pretrained path is replaced, residualized, bottlenecked, or unconstrained. The hard-token runs establish that relation signal is reachable but unsafe under replacement. The generic-negative run establishes that ordinary hard-negative mining does not reproduce the relation gain. The no-channel controls establish that a residual path without the bottleneck can still damage retrieval badly. These comparisons are the reason the method section emphasizes architecture and objective together.

\paragraph{Preservation-oriented adaptation.}
Residual adapters, low-rank adaptation, and null-space or principal-subspace constraints are natural tools for changing a pretrained model while limiting drift~\cite{hu2022lora,peng2025gnsp,luo2026keeplora,qiao2024pegp,wu2026psoft,lu2024vptnullspace}. We do not claim that residual adaptation or orthogonal projection is novel. In fact, the ROP follow-up in this project is framed as a preservation ablation and future direction, not as the headline method. The supported contribution is more specific: structured relational counterfactual training through a bounded residual interface, together with guardrail diagnostics showing where this interface helps and where it fails.

The missing standard PEFT baselines are an important limitation. LoRA, text-only tuning, vision-only tuning, prompt tuning, and related low-rank or subspace methods could offer different relation-preservation tradeoffs. This draft does not claim that RRB dominates them. Its narrower question is bottleneck-versus-unconstrained residual routing, not rank-constrained adaptation. Within the implemented run family, the bounded multi-channel residual interface avoids the most severe failures observed in hard replacement and no-channel residual controls. A stronger submission would need to place RRB against common adaptation baselines directly.
